\documentclass[sigconf,review]{acmart}
%%
%% \BibTeX command to typeset BibTeX logo in the docs
\AtBeginDocument{%
  \providecommand\BibTeX{{%
    Bib\TeX}}}

\usepackage[template]{mom-template} % remove "template" option to hide blue-font notes
\usepackage{multirow}
\usepackage{here}

\usepackage{xfp}
\makeatletter

\usepackage{xcolor}
\iftemplate
  \newcommand{\note}[1]{\textcolor{blue}{\textit{#1}}}
\else
  \newcommand{\note}[1]{}
\fi

% ============================================================
% CONFIGURATION
% ============================================================

% Total number of MTs
\def\MoM@mtcount{5}

% Number of questions per MT
\expandafter\def\csname MT@numq@1\endcsname{7}
\expandafter\def\csname MT@numq@2\endcsname{7}
\expandafter\def\csname MT@numq@3\endcsname{6}
\expandafter\def\csname MT@numq@4\endcsname{4}
\expandafter\def\csname MT@numq@5\endcsname{6}

% ============================================================
% INTERNALS -- everything below auto-derives from the above
% ============================================================

% \MTNumQ{n} -- expands to the number of questions for MT n
\newcommand{\MTNumQ}[1]{\csname MT@numq@#1\endcsname}

% \MT{<mt>}{<q>}{<score>}
% - if <score> is non-empty: store it
% - if <score> is empty: retrieve it (undefined entries read as 0)
\newcommand{\MT}[3]{%
  \if\relax\detokenize{#3}\relax
    \ifcsname MT@#1@#2\endcsname
      \csname MT@#1@#2\endcsname
    \else
      0%
    \fi
  \else
    \expandafter\gdef\csname MT@#1@#2\endcsname{#3}%
  \fi
}

% --- Loop counters (two to avoid conflicts in nested loops) ---
\newcount\MT@ctrA
\newcount\MT@ctrB

% --- Per-MT sum builder ---
% Builds an \fpeval-ready expression into \MT@sumexpr using \MT@ctrB.
\newcommand{\MT@BuildSum}[1]{%
  \global\MT@ctrB=0\relax
  \gdef\MT@sumexpr{0}%
  \loop
    \global\advance\MT@ctrB by 1
    \ifnum\MT@ctrB<\numexpr\MTNumQ{#1}+1\relax
      \xdef\MT@sumexpr{%
        \unexpanded\expandafter{\MT@sumexpr}%
        +\noexpand\MT{#1}{\the\MT@ctrB}{}}%
  \repeat
}

% \MTT{n} -- total score for MT n
\newcommand{\MTT}[1]{%
  \MT@BuildSum{#1}%
  \fpeval{\MT@sumexpr}%
}

% --- Grand sum builder ---
% Pre-computes each included MT's sum and wraps them in nested \fpeval calls
% so the final \fpeval in \MTG only sees numbers and +.
\newcommand{\MT@MaybeAddMT}[1]{%
  \ifcsname MT@included@#1\endcsname
    \begingroup
    \MT@BuildSum{#1}%
    \endgroup
    \xdef\MT@grandsumexpr{%
      \unexpanded\expandafter{\MT@grandsumexpr}%
      +\noexpand\fpeval{\unexpanded\expandafter{\MT@sumexpr}}}%
  \fi
}

\newcommand{\MT@BuildGrandSum}{%
  \global\MT@ctrA=0\relax
  \gdef\MT@grandsumexpr{0}%
  \loop
    \global\advance\MT@ctrA by 1
    \ifnum\MT@ctrA<\numexpr\MoM@mtcount+1\relax
      \MT@MaybeAddMT{\the\MT@ctrA}%
  \repeat
}

% \MTG -- grand total (only MTs marked as included)
\newcommand{\MTG}{%
  \MT@BuildGrandSum
  \fpeval{\MT@grandsumexpr}%
}

% --- Table row builder ---
% Builds \MT@Row calls into \MT@builtrows using \MT@ctrB.
\newcommand{\MT@Row}[2]{%
  & \textbf{#2} & \textsf{\MT{#1}{#2}{}} & \\
}

\newcommand{\MT@BuildRows}[1]{%
  \global\MT@ctrB=1\relax
  \gdef\MT@builtrows{}%
  \loop
    \global\advance\MT@ctrB by 1
    \ifnum\MT@ctrB<\numexpr\MTNumQ{#1}+1\relax
      \xdef\MT@builtrows{%
        \unexpanded\expandafter{\MT@builtrows}%
        \noexpand\MT@Row{#1}{\the\MT@ctrB}}%
  \repeat
}

% Full multirow block for MT #1
\newcommand{\MT@Block}[1]{%
  \MT@BuildRows{#1}%
  \multirow{\MTNumQ{#1}}{4em}{\textsf{\textbf{MT.#1}}}%
  & \textbf{1} & \textsf{\MT{#1}{1}{}} &
    \multirow{\MTNumQ{#1}}{*}{\textsf{\MTT{#1}}} \\
  \MT@builtrows
  \hline
}

% --- Summary table builder ---
% Builds \MT@IfIncluded calls into \MT@summaryblocks using \MT@ctrA.
\newcommand{\MT@IfIncluded}[1]{%
  \ifcsname MT@included@#1\endcsname\MT@Block{#1}\fi
}

\newcommand{\MT@BuildSummaryBlocks}{%
  \global\MT@ctrA=0\relax
  \gdef\MT@summaryblocks{}%
  \loop
    \global\advance\MT@ctrA by 1
    \ifnum\MT@ctrA<\numexpr\MoM@mtcount+1\relax
      \xdef\MT@summaryblocks{%
        \unexpanded\expandafter{\MT@summaryblocks}%
        \noexpand\MT@IfIncluded{\the\MT@ctrA}}%
  \repeat
}

% ============================================================
% PUBLIC INTERFACE
% ============================================================

% \MTScoreTable{n} -- insert a score table for MT n and mark it as included
\newcommand{\MTScoreTable}[1]{%
  \expandafter\gdef\csname MT@included@#1\endcsname{1}%
  \medskip
  \begin{center}
  \begin{tabular}{ |c|c|c|c| }
  \cline{2-4}
  \multicolumn{1}{c|}{} & \textbf{Question} & \textbf{Score} & \textbf{Total} \\
  \hline
  \MT@Block{#1}%
  \end{tabular}
  \end{center}
}

% \MTSummaryTable -- insert the appendix summary table (only included MTs)
\newcommand{\MTSummaryTable}{%
  \MT@BuildSummaryBlocks
  \begin{center}
  \begin{tabular}{ |c|c|c|c| }
  \cline{2-4}
  \multicolumn{1}{c|}{} & \textbf{Question} & \textbf{Score} & \textbf{Total} \\
  \hline
  \MT@summaryblocks
  \multicolumn{3}{|r|}{\textbf{Grand Total}} & \textsf{\MTG} \\
  \hline
  \end{tabular}
  \end{center}
}

\makeatother

% MT1 (7 questions)
\MT{1}{1}{0} 
\MT{1}{2}{0} 
\MT{1}{3}{0} 
\MT{1}{4}{0} 
\MT{1}{5}{0} 
\MT{1}{6}{0} 
\MT{1}{7}{0}

% MT2 (7 questions)
\MT{2}{1}{1} 
\MT{2}{2}{1} 
\MT{2}{3}{1} 
\MT{2}{4}{1} 
\MT{2}{5}{0.5} 
\MT{2}{6}{1} 
\MT{2}{7}{1}

% MT3 (6 questions)
\MT{3}{1}{0} 
\MT{3}{2}{0} 
\MT{3}{3}{0} 
\MT{3}{4}{1} 
\MT{3}{5}{1} 
\MT{3}{6}{0}

% MT4 (4 questions)
\MT{4}{1}{0} 
\MT{4}{2}{0} 
\MT{4}{3}{0} 
\MT{4}{4}{0}

% MT5 (6 questions)
\MT{5}{1}{0} 
\MT{5}{2}{0} 
\MT{5}{3}{0} 
\MT{5}{4}{0} 
\MT{5}{5}{0} 
\MT{5}{6}{0}

%% Rights management information.  This information is sent to you
%% when you complete the rights form.  These commands have SAMPLE
%% values in them; it is your responsibility as an author to replace
%% the commands and values with those provided to you when you
%% complete the rights form.
\setcopyright{acmlicensed}
\copyrightyear{2018}
\acmYear{2018}
\acmDOI{XXXXXXX.XXXXXXX}
%% These commands are for a PROCEEDINGS abstract or paper.
\acmConference[Conference acronym 'XX]{Make sure to enter the correct
  conference title from your rights confirmation email}{June 03--05,
  2018}{Woodstock, NY}
%%
%%  Uncomment \acmBooktitle if the title of the proceedings is different
%%  from ``Proceedings of ...''!
%%
%%\acmBooktitle{Woodstock '18: ACM Symposium on Neural Gaze Detection,
%%  June 03--05, 2018, Woodstock, NY}
\acmISBN{978-1-4503-XXXX-X/2018/06}


%%
%% Submission ID.
%% Use this when submitting an article to a sponsored event. You'll
%% receive a unique submission ID from the organizers
%% of the event, and this ID should be used as the parameter to this command.
%%\acmSubmissionID{123-A56-BU3}

%%
%% For managing citations, it is recommended to use bibliography
%% files in BibTeX format.
%%
%% You can then either use BibTeX with the ACM-Reference-Format style,
%% or BibLaTeX with the acmnumeric or acmauthoryear sytles, that include
%% support for advanced citation of software artefact from the
%% biblatex-software package, also separately available on CTAN.
%%
%% Look at the sample-*-biblatex.tex files for templates showcasing
%% the biblatex styles.
%%

%%
%% The majority of ACM publications use numbered citations and
%% references.  The command \citestyle{authoryear} switches to the
%% "author year" style.
%%
%% If you are preparing content for an event
%% sponsored by ACM SIGGRAPH, you must use the "author year" style of
%% citations and references.  Uncommenting
%% the next command will enable that style.
%%\citestyle{acmauthoryear}

%%
%% end of the preamble, start of the body of the document source.
\begin{document}

%%
%% The "title" command has an optional parameter,
%% allowing the author to define a "short title" to be used in page headers.
\title{MoM Challenge Response}

%%
%% The "author" command and its associated commands are used to define
%% the authors and their affiliations.
%% Of note is the shared affiliation of the first two authors, and the
%% "authornote" and "authornotemark" commands
%% used to denote shared contribution to the research.
\author{Jane Doe}
\email{author@email.com}
\orcid{1234-5678-9012}
\affiliation{%
  \institution{Institute}
  \city{City}
  \state{State}
  \country{Country}
}

\author{Niels Bohr, Max Planck}
\email{authors@same.institution}
\orcid{1234-5678-9012}
\affiliation{%
  \institution{Institute}
  \city{City}
  \state{State}
  \country{Country}
}



%%
%% By default, the full list of authors will be used in the page
%% headers. Often, this list is too long, and will overlap
%% other information printed in the page headers. This command allows
%% the author to define a more concise list
%% of authors' names for this purpose.
\renewcommand{\shortauthors}{Doe, J. \emph{et al.}}

%%
%% The abstract is a short summary of the work to be presented in the
%% article.
\begin{abstract}
We contribute a solution to the Model Management (MoM) Challenge 2026 
(available at doi: 10.5281/zenodo.19608752) with the XYZ tool / framework. 

\textcolor{red}{Describe how your contribution addressed the challenge 
requirements on a high abstraction level}

\end{abstract}

\maketitle

\section{Introduction}
\label{sec:Introduction}
\textcolor{red}{
General introduction of the solution:
\begin{itemize}
    \item Why did you chose to take on the Challenge?
    \item What are the general capabilities of your tool?
    \item Since when is the tool being developed?
    \item Is it open-source? 
    \item Is it dependent on another platform (Eclipse, etc.)
    \item Is it a commercial tool? Who are the customers?
\end{itemize}
}
\section{High-Level Dimensions and Functions of MoM}
\label{sec:HDF}

\textcolor{red}{Position your tool in the context of high-level dimensions 
and functions of MoM as described in the Challenge.}

\subsection{Approach [HDF.1]}
\label{subsec:HDF-1}

\subsection{Technological Space [HDF.2]}
\label{subsec:HDF-2}

\subsection{Relationships [HDF.3]}
\label{subsec:HDF-3}

\subsection{Views [HDF.4]}
\label{subsec:HDF-4}

\subsection{Collaboration [HDF.5]}
\label{subsec:HDF-5}

\section{Model Management Tasks}
\label{sec:MT}

\note{To enable automatic filling and computing of the score tables 
in this Section and \autoref{app:Scores}, locate the 
\texttt{score-tables.tex} \LaTeX{} file and change the scores in the last 
argument of the \texttt{\textbackslash{}MT} macros.  For example, giving 
1 point to \textsf{MT.2.3} translates to 
\texttt{\textbackslash{}MT\{2\}\{3\}\{1\}}.
%
If an MT is not addressed by your solution, you may remove its subsection 
and omit the corresponding \texttt{\textbackslash{}MTScoreTable} call; 
it will be excluded from the summary table automatically.}


%%% --------------------------------------------------------
%%% MT.1  Change Impact & Reconciliation
%%% --------------------------------------------------------
\subsection{Change Impact \& Reconciliation [MT.1]}
\label{sec:MT1}

\subsubsection*{MT.1.1 --- How is the change in the Parts Catalogue detected?}
\label{sec:MT1-1}

\note{Score 1 if your tool can detect and report changes; 0 otherwise.}


\subsubsection*{MT.1.2 --- Which mechanisms or processes ensure propagation of this change to the Bill of Materials?}
\label{sec:MT1-2}

\note{Score 1 if a propagation mechanism exists; 0 otherwise.}


\subsubsection*{MT.1.3 --- Is the propagation automatic, semi-automatic, or manual?}
\label{sec:MT1-3}

\note{Score 1 if propagation is at least semi-automatic; 0.5 if manual 
but documented; 0 if unsupported.}


\subsubsection*{MT.1.4 --- Can the propagation strategy be parameterised (to happen semi-/automatically)?}
\label{sec:MT1-4}

\note{Score 1 if the strategy is configurable; 0 otherwise.}


\subsubsection*{MT.1.5 --- At which granularity does propagation operate (e.g., element, feature, model)?}
\label{sec:MT1-5}

\note{Score 1 if granularity can be described; 0 if propagation is 
unsupported.}


\subsubsection*{MT.1.6 --- Can an explicit structure (e.g., graph, tree) be computed to represent impacted elements?}
\label{sec:MT1-6}

\note{Score 1 if an impact structure can be produced; 0 otherwise.}


\subsubsection*{MT.1.7 --- Can such structures be composed across multiple changes, and how deeply can they be inspected?}
\label{sec:MT1-7}

\note{Score 1 if composition and inspection are supported; 0 otherwise.}


\subsubsection*{Alternatives \& Extensions}
\label{sec:MT1-Alternatives-Extensions}

\textcolor{red}{Discuss alternative mechanisms or extensions beyond the 
scope of the questions above, if any.}

\subsubsection*{Score}
\label{sec:MT1-Score}

\MTScoreTable{1}


%%% --------------------------------------------------------
%%% MT.2  Views Management
%%% --------------------------------------------------------
\subsection{Views Management [MT.2]}
\label{sec:MT2}

\subsubsection*{MT.2.1 --- Can model elements or types be selectively filtered on demand? On which criteria?}
\label{sec:MT2-1}

\note{Score 1 if selective filtering is supported; 0 otherwise.}


\subsubsection*{MT.2.2 --- Can filtering rules be specified declaratively (e.g., via a query language)?}
\label{sec:MT2-2}

\note{Score 1 if declarative specification is supported; 0 otherwise.}


\subsubsection*{MT.2.3 --- Can views be systematically derived from conformance models? How?}
\label{sec:MT2-3}

\note{Score 1 if systematic derivation is supported; 0 otherwise.}


\subsubsection*{MT.2.4 --- Are views themselves typed (e.g., via a metamodel)?}
\label{sec:MT2-4}

\note{Score 1 if views are typed; 0 otherwise.}


\subsubsection*{MT.2.5 --- How are changes in views reconciled with the source model?}
\label{sec:MT2-5}

\note{Score 1 if view-to-source reconciliation is supported; 0 otherwise.}


\subsubsection*{MT.2.6 --- How are views notified of changes in the source model?}
\label{sec:MT2-6}

\note{Score 1 if notification or synchronisation mechanism exists; 
0 otherwise.}


\subsubsection*{MT.2.7 --- Do source model changes automatically propagate to views? How?}
\label{sec:MT2-7}

\note{Score 1 if automatic propagation is supported; 0 otherwise.}


\subsubsection*{Alternatives \& Extensions}
\label{sec:MT2-Alternatives-Extensions}

\textcolor{red}{Discuss alternative mechanisms or extensions beyond the 
scope of the questions above, if any.}

\subsubsection*{Score}
\label{sec:MT2-Score}

\MTScoreTable{2}


%%% --------------------------------------------------------
%%% MT.3  Querying & Validation
%%% --------------------------------------------------------
\subsection{Querying \& Validation [MT.3]}
\label{sec:MT3}

\subsubsection*{MT.3.1 --- Is \textsf{calculated\_total\_mass\_kg} automatically updated after the change?}
\label{sec:MT3-1}

\note{Score 1 if automatic update is supported; 0 otherwise.}


\subsubsection*{MT.3.2 --- Which mechanisms or processes support this update?}
\label{sec:MT3-2}

\note{Score 1 if a concrete mechanism is described; 0 otherwise.}


\subsubsection*{MT.3.3 --- Does it require manual intervention?}
\label{sec:MT3-3}

\note{Score 1 if no manual intervention is needed; 0.5 if partially 
manual; 0 if fully manual or unsupported.}


\subsubsection*{MT.3.4 --- Can the process be fully automated or parameterised?}
\label{sec:MT3-4}

\note{Score 1 if automation is configurable; 0 otherwise.}


\subsubsection*{MT.3.5 --- Is the updated status reflected in the Report?}
\label{sec:MT3-5}

\note{Score 1 if the Report is updated accordingly; 0 otherwise.}


\subsubsection*{MT.3.6 --- Can impact analysis be performed manually by querying all relevant models?}
\label{sec:MT3-6}

\note{Score 1 if cross-model querying is supported; 0 otherwise.}


\subsubsection*{Alternatives \& Extensions}
\label{sec:MT3-Alternatives-Extensions}

\textcolor{red}{Discuss alternative mechanisms or extensions beyond the 
scope of the questions above, if any.}

\subsubsection*{Score}
\label{sec:MT3-Score}

\MTScoreTable{3}


%%% --------------------------------------------------------
%%% MT.4  Concurrent Modifications
%%% --------------------------------------------------------
\subsection{Concurrent Modifications [MT.4]}
\label{sec:MT4}

\subsubsection*{MT.4.1 --- How are conflicts resolved: manual vs.\ (semi-)automatic, online vs.\ offline, configurable vs.\ fixed strategies?}
\label{sec:MT4-1}

\note{Score 1 if a conflict resolution strategy is supported; 
0 otherwise.}


\subsubsection*{MT.4.2 --- Are changes propagated in real time or through explicit synchronization?}
\label{sec:MT4-2}

\note{Score 1 if propagation mode is described; 0 if concurrent 
editing is unsupported.}


\subsubsection*{MT.4.3 --- If concurrent conflicting modifications are prevented, which mechanisms or policies detect and avoid such conflicts?}
\label{sec:MT4-3}

\note{Score 1 if prevention mechanisms are described; 0 if neither 
prevention nor resolution is supported.}


\subsubsection*{MT.4.4 --- Does the tool support live collaborative modelling? If so, how are conflicts visualized, prevented, or resolved?}
\label{sec:MT4-4}

\note{Score 1 if live collaboration is supported; 0 otherwise.}


\subsubsection*{Alternatives \& Extensions}
\label{sec:MT4-Alternatives-Extensions}

\textcolor{red}{Discuss alternative mechanisms or extensions beyond the 
scope of the questions above, if any.}

\subsubsection*{Score}
\label{sec:MT4-Score}

\MTScoreTable{4}


%%% --------------------------------------------------------
%%% MT.5  Version Management
%%% --------------------------------------------------------
\subsection{Version Management [MT.5]}
\label{sec:MT5}

\subsubsection*{MT.5.1 --- What is the definition of a ``version'' (single artefact, dependency closure, or system-wide state)?}
\label{sec:MT5-1}

\note{Score 1 if the notion of version is clearly defined; 0 otherwise.}


\subsubsection*{MT.5.2 --- Are versions created automatically or explicitly (e.g., commit-based)?}
\label{sec:MT5-2}

\note{Score 1 if version creation is supported; 0 otherwise.}


\subsubsection*{MT.5.3 --- How are artefacts linked across versions (e.g., traceability links between instances)?}
\label{sec:MT5-3}

\note{Score 1 if cross-version linking is supported; 0 otherwise.}


\subsubsection*{MT.5.4 --- Can versions be enriched with metadata?}
\label{sec:MT5-4}

\note{Score 1 if metadata attachment is supported; 0 otherwise.}


\subsubsection*{MT.5.5 --- Are such metadata queryable and exploitable?}
\label{sec:MT5-5}

\note{Score 1 if metadata can be queried; 0 otherwise.}


\subsubsection*{MT.5.6 --- Can the rationale for version creation be captured and maintained?}
\label{sec:MT5-6}

\note{Score 1 if rationale capture is supported; 0 otherwise.}


\subsubsection*{Alternatives \& Extensions}
\label{sec:MT5-Alternatives-Extensions}

\textcolor{red}{Discuss alternative mechanisms or extensions beyond the 
scope of the questions above, if any.}

\subsubsection*{Score}
\label{sec:MT5-Score}

\MTScoreTable{5}

\section{Replication Package (Optional)}
\label{sec:RP}

\section{Conclusion}
\label{sec:Conclusion}

\textcolor{red}{
\begin{itemize}
   \item Summarize the \textbf{main strengths and limitations} of the approach;
   \item Describe any \textbf{adaptations or assumptions} made to fit the Case Study;
   \item Identify \textbf{unaddressed requirements or limitations} of the challenge;
   \item Provide \textbf{lessons learnt} and implications for future work,
   both for the authors and the broader MoM community.
\end{itemize}
}

\bibliographystyle{ACM-Reference-Format}
\bibliography{citations}

\appendix
\section{Scores Summary Table}
\label{app:Scores}

\MTSummaryTable

\end{document}
\endinput
%%
%% End of file `sample-sigconf-authordraft.tex'.
